\documentclass[final]{beamer}
\usepackage[size=a0,scale=1.0,orientation=portrait]{beamerposter}
\usetheme{default}

\usepackage[T1]{fontenc}
\usepackage[utf8]{inputenc}
\usepackage{amsmath,amssymb}
\usepackage{booktabs}
\usepackage{graphicx}
\usepackage{xcolor}
\usepackage{caption}

% Figures: works whether the zip was extracted as a folder or the PNGs were uploaded flat
\graphicspath{{Exp008_CNN_Transformer_Results/}{./}}

% ================= PRECISE COLOR PALETTE =================
\definecolor{IITBBlue}{RGB}{0, 43, 92}
\definecolor{AccentGold}{RGB}{196, 158, 54}
\definecolor{TextBlack}{RGB}{33, 37, 41}
\definecolor{DividerGray}{RGB}{210, 210, 210}

% ================= TYPOGRAPHY =================
\usepackage{helvet}
\renewcommand{\familydefault}{\sfdefault}
\setbeamercolor{background canvas}{bg=white}
\setbeamertemplate{navigation symbols}{}

% ================= BLOCK DESIGN (Fixed Double Title) =================
\setbeamercolor{block title}{fg=IITBBlue, bg=white}
\setbeamercolor{block body}{fg=TextBlack, bg=white}
\setbeamerfont{block title}{size=\huge, series=\bfseries}
\setbeamerfont{block body}{size=\large}

\setbeamertemplate{block begin}{
    \vspace{0.8em}
    \begin{beamercolorbox}[leftskip=0ex]{block title}
        {\color{AccentGold}\rule[0.2em]{1.5cm}{6pt}} \hspace{0.2em} \insertblocktitle
    \end{beamercolorbox}
    \vspace{0.2em}
}
\setbeamertemplate{block end}{\vspace{0.4em}}

% ================= HEADER =================
\setbeamertemplate{headline}{
    \leavevmode

    \begin{beamercolorbox}[
        wd=\paperwidth,
        sep=0pt
    ]{headline}

    \vspace{1.2cm}

    \begin{columns}[c,totalwidth=\paperwidth]

        % LEFT LOGO
        \begin{column}{0.16\paperwidth}
            \centering
            \includegraphics[
                width=0.115\paperwidth,
                height=7cm,
                keepaspectratio
            ]{iitblogo.png}
        \end{column}

        % TITLE / AUTHORS
        \begin{column}{0.68\paperwidth}
            \centering

            {\color{IITBBlue}
            \fontsize{48}{56}
            \selectfont
            \bfseries
            GeoDrillNet: Deep Learning for Continuous\\
            Wellbore Geology Prediction from Drilling Measurements
            \par}

            \vspace{0.45cm}

            {\color{TextBlack}
            \fontsize{27}{34}
            \selectfont
            Priyanshu Sharma$^{1}$ \quad
            Rahul Sharma$^{2}$ \quad
            Aakash R. Nair$^{3}$ \quad
            Bharath Shekar$^{3}$
            \par}

            \vspace{0.3cm}

            {\color{AccentGold}
            \fontsize{20}{26}
            \selectfont
            $^{1}$IIT Madras \quad
            $^{2}$Kurukshetra University \quad
            $^{3}$IIT Bombay
            \par}

        \end{column}

        % RIGHT SIMLAB LOGO
        \begin{column}{0.16\paperwidth}
            \centering
            \includegraphics[
                width=0.13\paperwidth,
                height=5.5cm,
                keepaspectratio
            ]{simlab2.png}
        \end{column}

    \end{columns}

    \vspace{0.5cm}

    {\centering
        {\color{DividerGray}
        \rule{0.92\paperwidth}{2.5pt}}\par}

    \vspace{0.1cm}

    \end{beamercolorbox}
}

% ================= LAYOUT =================
\newlength{\sepwidth}
\newlength{\colwidth}
\setlength{\sepwidth}{0.03\paperwidth}
\setlength{\colwidth}{0.45\paperwidth}

\begin{document}
\begin{frame}
\vspace{-3.2cm}
\begin{columns}[T,totalwidth=\paperwidth]
\hspace{\sepwidth}

% ================= COLUMN 1 =================
\begin{column}{\colwidth}
    \begin{block}{Abstract}
        Continuous wellbore geology prediction is fundamental to automated geosteering during horizontal
        drilling, yet existing data-driven approaches often struggle to capture the complex spatial and
        temporal relationships present in drilling measurements. We propose \textbf{GeoDrillNet}, a deep
        learning framework that learns geological representations directly from sequential drilling data.
        GeoDrillNet combines hierarchical convolutional feature extraction with Transformer self-attention
        to characterize both local geological variations and long-range formation trends, predicting the
        geological position (TVT) along the well trajectory. Out-of-fold evaluation on 773 wells of the
        ROGII Wellbore Geology Prediction benchmark gives an RMSE of \textbf{21.33\,ft} on the prediction
        interval, matching a larger TCN--cross-attention variant (21.40\,ft) with
        \textbf{5$\times$ fewer parameters}.
    \end{block}

    \begin{block}{Motivation \& Problem}
        \begin{itemize}
            \item Geosteering keeps a horizontal well inside the target formation; the geological
                  position (TVT) is picked manually by correlating gamma ray (GR) with a vertical typewell
                  --- slow, subjective, and hard to scale to real time.
            \item \textbf{Target:} predict TVT beyond the \textbf{Prediction Start (PS)} point from the
                  well path (X, Y, Z), GR, and the typewell log.
            \item \textbf{Metric:} $\mathrm{RMSE} = \sqrt{\tfrac{1}{N}\sum_i
                  (\mathrm{TVT}^{\mathrm{manual}}_i - \mathrm{TVT}^{\mathrm{pred}}_i)^2}$ over predicted points.
        \end{itemize}
    \end{block}

    \begin{block}{Dataset}
        773 horizontal wells (one-foot sampling) with an assigned typewell each. Wells are split by
        5-fold GroupKFold so every well is evaluated by a model that never saw it; colour shows each
        well's out-of-fold error.
        \par\vspace{0.4cm}
        \centering
        \includegraphics[width=0.7\linewidth]{06_map_view_error.png}
    \end{block}

    \begin{block}{Proposed: GeoDrillNet (CNN + Transformer)}
        Residual 1D-CNN blocks capture local log character; a Transformer encoder relates every foot of
        the window to the whole lateral. Target standardized per fold; overlapping windows and 5 fold
        models are averaged.
        \par\vspace{0.4cm}
        \centering
        \includegraphics[width=0.95\linewidth]{01_architecture_diagram.png}
    \end{block}

    \begin{block}{Conclusions}
        \begin{itemize}
            \item \textbf{Efficiency:} 0.37M parameters, 5$\times$ fewer than the TCN--cross-attention
                  model, at the same accuracy after PS (21.33 vs.\ 21.40\,ft).
            \item \textbf{Accuracy:} median per-well RMSE of 14.5\,ft; the best held-out wells stay within
                  9--12\,ft of the manual TVT over 2,500--5,500\,ft of prediction.
            \item \textbf{Limitation:} error grows with distance from PS; a constant-TVT baseline
                  (15.91\,ft) is still lower on the prediction interval.
            \item \textbf{Next steps:} predict $\Delta$TVT relative to the last known value, explicit
                  GR--typewell alignment, and stable mixed-precision training.
        \end{itemize}
    \end{block}

    \begin{block}{Hyper-Parameter Search}
        150-trial Optuna search (TPE sampler, median pruning) over CNN/Transformer size, dropout,
        learning rate and window length.
        \par\vspace{0.3cm}
        \centering
        \includegraphics[width=0.62\linewidth]{17_optuna_search_history.png}
    \end{block}
\end{column}

\hspace{\sepwidth}

% ================= COLUMN 2 =================
\begin{column}{\colwidth}
    \begin{block}{Model Benchmarking}
        \centering
        \renewcommand{\arraystretch}{1.5}
        \resizebox{\linewidth}{!}{%
        \begin{tabular}{l c c c}
            \toprule
            \textbf{Model} & \textbf{Params} & \textbf{Train Time} & \textbf{RMSE after PS} \\
            \midrule
            \textbf{GeoDrillNet (CNN + Transformer)} & \textbf{0.37M} & \textbf{684 s} & \textbf{21.33 ft} \\
            TCN + Cross-Attention + Transformer & 1.91M & 1238 s & 21.40 ft \\
            Hold last known TVT constant & -- & -- & 15.91 ft \\
            \bottomrule
        \end{tabular}}

        \vspace{0.3cm}
        {\normalsize Out-of-fold, 5-fold GroupKFold by well; train time = 5-fold CV on one GPU.
        A 51-ft moving average of GeoDrillNet's output gives 21.16\,ft.}
    \end{block}

    \begin{block}{Training Dynamics}
        \centering
        \includegraphics[width=0.95\linewidth]{02_loss_curves.png}
    \end{block}

    \begin{block}{Prediction Quality (Out-of-Fold)}
        \centering
        \includegraphics[width=0.95\linewidth]{04_prediction_quality.png}
    \end{block}

    \begin{block}{Held-Out Well Predictions}
        Wells never seen in training, predicted beyond PS over 2,500--5,500\,ft: GeoDrillNet stays within
        9--12\,ft of the manual interpretation along the entire lateral (RMSE 3.2--5.7\,ft).
        \par\vspace{0.4cm}
        \centering
        \includegraphics[width=0.95\linewidth]{22_heldout_predictions_low_error.png}
    \end{block}

    \begin{block}{Per-Well Error Distribution}
        \centering
        \includegraphics[width=0.95\linewidth]{05_per_well_rmse.png}
    \end{block}
\end{column}

\hspace{\sepwidth}
\end{columns}

\end{frame}
\end{document}
